% !TEX root = ../../../main.tex
% Block card: detection.wavelet_scattering -- "Wavelet scattering" (Part I, the blocks chapter).
\cardfield{What it does}
Summarises, stretch by stretch, how the signal's energy is spread over slow and fast rhythms,
in a form that changes little if an event shifts slightly in time.

\cardfield{When to use it}
When you want to compare or group stretches of noisy recordings by their texture rather than
look for individual events: a set of features, not a detector.

\cardfield{What goes in, what comes out}
In goes a span of one channel. Out comes an image: one row per feature (each a band of rhythm,
or how that band's strength itself rises and falls), one column per block of $2^J$ samples
(256\,s at 1\,Hz with the default \emph{J}).

\cardfield{Settings that matter}
\emph{J} sets the longest timescale and also the length of each output step, $2^J$ samples;
the span must hold at least that many. \emph{Q} is how finely each doubling of timescale is
divided. \emph{max\_order} 2 adds the ``how the strength varies'' features; 1 leaves only the
bands.

\cardfield{Watch out for}
Its columns are blocks of samples, not samples, so it cannot feed
\blockref{detection.wavelet_summation}. On a span not much longer than $2^J$ samples the ends
are distorted by padding; the block notes this as \code{border_effects_flagged}.
